\subsection{Результаты в среде CartPole}

Первым этапом экспериментального исследования стала классическая задача управления — балансировка перевернутого маятника (CartPole). Данная среда выступала в качестве фундаментального полигона для верификации работоспособности алгоритмов и оценки базовых характеристик когнитивных агентов. В рамках эксперимента проводилось сопоставление классического алгоритма глубокого Q-обучения (DQN) и нейроэволюционного метода (NEAT).

На этапе проектирования архитектур моделей были применены два принципиально разных подхода. Агент на базе DQN использовал фиксированную полносвязную нейронную сеть с одним скрытым слоем, содержащую 114 параметров. В противовес этому, алгоритм NEAT функционировал в парадигме постепенного усложнения (complexifying), начиная эволюционный процесс с простейшей топологии без скрытых слоев и динамически формируя структуру сети в ответ на требования среды.

Сравнительный анализ метрик (согласно определениям \ref{def:learning_efficiency}-\ref{def:robustness}) позволил сделать следующие ключевые выводы:


\subsubsection{Эффективность использования выборки и скорость схождения}

Оба алгоритма успешно справились с задачей, однако нейроэволюционный подход продемонстрировал существенно более высокую скорость адаптации. Алгоритму NEAT потребовалось всего 4 поколения (или 10 548 шагов взаимодействия со средой), чтобы найти оптимальную стратегию. В то же время алгоритму DQN для стабилизации весов и схождения потребовалось 471 эпизод, что эквивалентно 25 923 взаимодействиям. Таким образом, эволюционный алгоритм оказался более чем в 2,4 раза эффективнее с точки зрения затрат вычислительных тактов симуляции.


\subsubsection{Качество решения и архитектурная сложность}

В идеальных условиях обе модели достигли абсолютного максимума качества (средняя награда 500.00 из 500) со 100\% стабильностью (стандартное отклонение равно 0). Однако способы достижения этого результата кардинально различаются. В то время как DQN использует плотную полносвязную структуру (плотность связей 100\%), NEAT синтезировал высокоэффективную разреженную сеть, состоящую всего из 9 параметров (плотность связей 75\%). Это доказывает способность нейроэволюции находить экстремально компактные и легковесные решения, что критически важно для развертывания моделей на устройствах с ограниченными вычислительными ресурсами.


\subsubsection{Устойчивость к внешним возмущениям}

Наиболее значимое различие между алгоритмами было зафиксировано при стресс-тестировании агентов в условиях зашумленных данных. Для генерации шума в этом и последующий проектах используется Гауссовский шум. В данной среде берется стандартное отклонение равное 0.1. При добавлении шума в наблюдения среды, эффективность полносвязной модели DQN резко снизилась на 29,14\% (средняя награда упала до 354.30). Агент оказался склонен к переобучению под «идеализированную» картину мира. Напротив, агент NEAT продемонстрировал абсолютную робастность: в условиях аналогичного зашумления модель не потеряла ни одного балла, сохранив среднюю награду на уровне 500.00. Столь высокая устойчивость напрямую вытекает из компактности эволюционной архитектуры: из-за отсутствия «лишних» связей и параметров, сеть физически лишена способности пропускать через себя шумовые сигналы, фильтруя их на уровне самой топологии.


\subsubsection{Практические выводы}

Эксперимент в среде CartPole эмпирически доказал, что для решения базовых задач когнитивного управления нейроэволюционный алгоритм NEAT превосходит классический DQN по скорости схождения, компактности получаемой модели и, что наиболее важно, по уровню устойчивости к неопределенностям среды.